%%%%%%%%%%%%%%%%%%%%%%%%%%%%%%%%%%%%%%%%%%%%%%%%%%%%%%%%%%%%%%%%%%%%%%%%%%%%%%%
\subsubsection{Domain 10.9.2: Health and Prevention}
\label{subsubsec:health-prevention}
%%%%%%%%%%%%%%%%%%%%%%%%%%%%%%%%%%%%%%%%%%%%%%%%%%%%%%%%%%%%%%%%%%%%%%%%%%%%%%%

\paragraph{The Challenge.}
Preventive health behaviors---quitting smoking, exercise, healthy diet---offer
large long-term benefits but require short-term sacrifices. Despite knowing
the consequences, people systematically underinvest in prevention. The EBF
framework explains the intertemporal structure of this failure.

\paragraph{Context Profile ($\Psi$).}

\begin{table}[htbp]
\centering
\caption{Health Behavior Context Factors}
\label{tab:health-psi}
\renewcommand{\arraystretch}{1.3}
\begin{tabular}{@{}lll@{}}
\toprule
\textbf{$\Psi$-Dimension} & \textbf{Effect} & \textbf{Mechanism} \\
\midrule
$\Psi_T$ (Temporal stability) & Critical & Low $\Psi_T$ → high discounting → ignore future \\
$\Psi_C$ (Cognitive) & Important & Health literacy enables informed decisions \\
$\Psi_S$ (Social trust) & Supportive & Social support enables behavior change \\
$\Psi_I$ (Institutional) & Enabling & Healthcare access, health promotion \\
$\Psi_E$ (Economic) & Ambiguous & Wealth enables but also enables indulgence \\
\bottomrule
\end{tabular}
\end{table}

\paragraph{Relevant Utility Components.}
Health behaviors involve a classic intertemporal trade-off:

\begin{table}[htbp]
\centering
\caption{Health Behavior: Utility Components by Time}
\label{tab:health-components}
\renewcommand{\arraystretch}{1.2}
\begin{tabular}{@{}lcccc@{}}
\toprule
\textbf{Component} & $t=0$ & $t=1$ & $t=2$ & $t=3$ \\
& Instant & Short & Medium & Long \\
\midrule
\multicolumn{5}{l}{\textit{Smoking Cessation Example}} \\
$U^{INU}_E$ (Emotional) & \textbf{--0.6} & --0.3 & +0.1 & +0.2 \\
$U^{INU}_P$ (Physical) & --0.2 & --0.1 & +0.3 & \textbf{+0.8} \\
$U^{INU}_F$ (Financial) & 0 & +0.1 & +0.2 & +0.3 \\
\midrule
\multicolumn{5}{l}{\textit{Exercise Adoption Example}} \\
$U^{INU}_E$ (Emotional) & \textbf{--0.4} & --0.2 & +0.2 & +0.3 \\
$U^{INU}_P$ (Physical) & --0.3 & 0 & +0.4 & \textbf{+0.7} \\
$U^{INU}_S$ (Social) & --0.1 & +0.1 & +0.2 & +0.2 \\
\bottomrule
\end{tabular}
\end{table}

\paragraph{The Core Conflict: $t=0$ vs. $t=3$.}
The fundamental problem is temporal:

\begin{align}
U^{INU}_{t=0} &< 0 \quad \text{(immediate pain: withdrawal, effort, discomfort)} \\
U^{INU}_{t=3} &> 0 \quad \text{(long-term gain: health, energy, longevity)}
\end{align}

With standard discounting ($\delta \approx 0.85$):
\begin{align}
\text{Present value of } U_{t=0} &= 1.0 \times (-0.6) = -0.60 \\
\text{Present value of } U_{t=3} &= (0.85)^3 \times (+0.8) = 0.61 \times 0.8 = +0.49
\end{align}

The discounted long-term benefit barely exceeds the immediate cost---and 
with any additional uncertainty or hyperbolic discounting, the behavior fails.

\paragraph{Reference Points in Health.}

\begin{table}[htbp]
\centering
\caption{Health Reference Points}
\label{tab:health-reference}
\renewcommand{\arraystretch}{1.3}
\begin{tabular}{@{}llp{5.5cm}@{}}
\toprule
\textbf{Category} & \textbf{Reference} & \textbf{Implication} \\
\midrule
INU (Status quo) & Current health state & Smoker's reference = ``smoker normal'' \\
INU (Expectation) & Expected health trajectory & Often unrealistically optimistic \\
IDN (Identity) & ``I'm a smoker'' / ``I'm not athletic'' & Identity blocks change \\
KNU (Group norm) & Peer health behaviors & Social smoking maintains behavior \\
\bottomrule
\end{tabular}
\end{table}

\paragraph{The Identity Dimension.}
Health behaviors often have strong IDN components:

\begin{itemize}
\item ``I'm a smoker'' --- smoking is part of self-concept
\item ``I'm not a gym person'' --- exercise violates identity
\item ``I'm a foodie'' --- diet restrictions threaten identity
\end{itemize}

Behavior change requires \emph{identity change}:
\begin{equation}
\Delta U^{IDN} = U^{IDN}_{new identity} - U^{IDN}_{old identity}
\end{equation}

If old identity is strongly held ($r^{IDN}$ sticky), $\Delta U^{IDN} < 0$ 
creates additional resistance.

\paragraph{The Social Dimension (KNU).}
Health behaviors are socially embedded:

\begin{itemize}
\item Smoking breaks are social rituals
\item Drinking is community bonding
\item Unhealthy food is cultural tradition
\end{itemize}

Changing behavior may mean:
\begin{equation}
\Delta U^{KNU}_S < 0 \quad \text{(leaving the smoking group, being ``difficult'' at dinners)}
\end{equation}

\paragraph{Welfare Calculation: Smoking Cessation.}
Full calculation including all dimensions:

\begin{table}[htbp]
\centering
\caption{Smoking Cessation: Complete Welfare Analysis}
\label{tab:smoking-welfare}
\renewcommand{\arraystretch}{1.2}
\begin{tabular}{@{}lrrr@{}}
\toprule
\textbf{Component} & \textbf{$\Delta U$} & \textbf{Weight} & \textbf{Contribution} \\
\midrule
\multicolumn{4}{l}{\textit{INU (discounted, $\delta = 0.85$)}} \\
$U^{INU}_{E,0}$ (withdrawal) & --0.60 & $\times 1.0$ & --0.60 \\
$U^{INU}_{P,3}$ (health) & +0.80 & $\times 0.61$ & +0.49 \\
$U^{INU}_{F,2}$ (savings) & +0.20 & $\times 0.72$ & +0.14 \\
\textit{INU subtotal} & & & \textbf{+0.03} \\
\midrule
\multicolumn{4}{l}{\textit{KNU}} \\
$U^{KNU}_{S,0}$ (leaving smoking group) & --0.30 & $\times 1.0$ & --0.30 \\
$U^{KNU}_{P,2}$ (family benefits) & +0.20 & $\times 0.72$ & +0.14 \\
\textit{KNU subtotal} & & & \textbf{--0.16} \\
\midrule
\multicolumn{4}{l}{\textit{IDN}} \\
$U^{IDN}_{E,0}$ (identity disruption) & --0.40 & $\times 1.0$ & --0.40 \\
$U^{IDN}_{E,3}$ (new identity) & +0.30 & $\times 0.61$ & +0.18 \\
\textit{IDN subtotal} & & & \textbf{--0.22} \\
\midrule
\multicolumn{3}{l}{\textbf{Total $\Delta Q$ (before weighting by $\omega$)}} & \textbf{--0.35} \\
\bottomrule
\end{tabular}
\end{table}

\textbf{Result:} Even with substantial long-term health benefits, 
the immediate INU pain, KNU disruption, and IDN violation produce 
\emph{negative} total welfare change. The ``irrational'' smoker is 
actually optimizing correctly given their welfare function.

\paragraph{Why Health Interventions Fail.}
Traditional interventions target only one component:

\begin{table}[htbp]
\centering
\caption{Why Health Interventions Fail}
\label{tab:health-failures}
\renewcommand{\arraystretch}{1.3}
\begin{tabular}{@{}lll@{}}
\toprule
\textbf{Intervention} & \textbf{Targets} & \textbf{Ignores} \\
\midrule
Information campaigns & $U^{INU}_{P,3}$ & IDN, KNU, $t=0$ costs \\
Financial incentives & $U^{INU}_{F}$ & IDN, KNU, emotional costs \\
Fear appeals & $P^{INU}_{P,3}$ (increase pain) & Backfires via reactance \\
Gym memberships & $\Delta U^{INU}_P$ & IDN (``not a gym person'') \\
\bottomrule
\end{tabular}
\end{table}

\paragraph{Effective Intervention Design.}

\begin{table}[htbp]
\centering
\caption{EBF Health Interventions}
\label{tab:health-interventions}
\renewcommand{\arraystretch}{1.3}
\begin{tabular}{@{}llp{5cm}@{}}
\toprule
\textbf{Target} & \textbf{Strategy} & \textbf{Example} \\
\midrule
$\delta$ (discounting) & Increase patience & Commitment devices, pre-commitment \\
$a_{INU,t=3}$ & Make future salient & Aging apps, health projections \\
$U^{IDN}$ & Identity transformation & ``Becoming a non-smoker'' programs \\
$U^{KNU}$ & Social transition & Group quitting, peer support \\
$r^{INU}_{P}$ & Shift reference point & ``Your lungs of a 60-year-old'' \\
$U^{INU}_{E,0}$ & Reduce immediate pain & Nicotine replacement, gradual change \\
\bottomrule
\end{tabular}
\end{table}

\paragraph{Case Example: Successful Smoking Cessation Program.}
The most effective programs combine multiple components:

\begin{enumerate}
\item \textbf{Pharmacological support:} Reduces $P^{INU}_{E,0}$ (withdrawal pain)
\item \textbf{Group therapy:} Maintains $U^{KNU}$ (new social support)
\item \textbf{Identity work:} Builds $U^{IDN}$ (``I am a non-smoker'')
\item \textbf{Commitment contracts:} Increases effective $\delta$
\item \textbf{Milestone rewards:} Creates $G^{INU}_{E,t}$ along the way
\end{enumerate}

Success rates increase from ~5\% (cold turkey) to ~30\% (comprehensive program)
because the intervention addresses the \emph{full welfare function}, not just 
health information.

\paragraph{Key Insight.}
\begin{tcolorbox}[colback=yellow!10!white, colframe=yellow!50!black]
\textbf{Health Behavior Insight:}
Prevention fails not because people don't know the facts, but because:
\begin{enumerate}
\item Future benefits are heavily discounted ($\delta^3 \approx 0.6$)
\item Immediate costs are fully felt ($\delta^0 = 1.0$)
\item Identity (``smoker'') blocks change
\item Social ties (``smoking group'') resist change
\end{enumerate}
Effective interventions must address \emph{all four} barriers simultaneously,
not just provide information about long-term health consequences.
\end{tcolorbox}

%%%%%%%%%%%%%%%%%%%%%%%%%%%%%%%%%%%%%%%%%%%%%%%%%%%%%%%%%%%%%%%%%%%%%%%%%%%%%%%
% END OF 10.9.2
%%%%%%%%%%%%%%%%%%%%%%%%%%%%%%%%%%%%%%%%%%%%%%%%%%%%%%%%%%%%%%%%%%%%%%%%%%%%%%%
